% C.3ax. Complete root proof and independent Frobenius audit; restricted cone result.
\paragraph{C.3ax (Proved): rational rigidity for smooth projective cones in positive characteristic.}
Let $k$ be an algebraically closed field of characteristic $p>0$,
let $V/k$ be a smooth connected projective variety of dimension $d\ge1$,
and let $L$ be ample. Suppose the full section ring
$S=\bigoplus_{n\ge0}H^0(V,L^n)$ is normal and generated in degree one.
Let $R=(S_{S_+})^{\mathrm{sh}}$, let
$X=\operatorname{Spec}R\setminus\{\mathfrak m_R\}$, and let
$Z=\operatorname{Tot}(L^{-1})\times_{\operatorname{Spec}S}\operatorname{Spec}R$.
Write $i:V\hookrightarrow Z$ for the zero section and $p_V:Z\to V$
for the projection. Then restriction and pullback give mutually inverse maps
\begin{equation}\label{eq:C3ax-K}
 i^*:K_0(Z)\otimes\mathbb Z[1/p]
       \ \xrightarrow{\ \sim\ }\ K_0(V)\otimes\mathbb Z[1/p],
 \qquad (i^*)^{-1}=p_V^*.
\end{equation}
For every $r\ge0$ we consequently have
\begin{align}
 i^!: \mathrm{CH}^r(Z)_{\mathbb Q}
     &\xrightarrow{\ \sim\ }\mathrm{CH}^r(V)_{\mathbb Q},
          & (i^!)^{-1}&=p_V^*,\label{eq:C3ax-Chow}\\
 \mathrm{CH}^r(X)_{\mathbb Q}
     &\simeq\frac{\mathrm{CH}^r(V)_{\mathbb Q}}
                       {c_1(L)\mathrm{CH}^{r-1}(V)_{\mathbb Q}},\label{eq:C3ax-open}\\
 B^r(X)_{\mathbb Q}
     &\simeq\frac{B^r(V)_{\mathbb Q}}
                       {c_1(L)B^{r-1}(V)_{\mathbb Q}}.\label{eq:C3ax-B}
\end{align}
Here $B$ is the quotient by ordinary algebraic equivalence from actual
fixed-product families. In particular $B^d(X)_{\mathbb Q}=0$.
The assertion concerns these cone germs. It does not prove rational
top-degree finiteness for arbitrary strictly local singularities.

\begin{proof}
\emph{1. Finite detection without a characteristic-zero hypothesis.}
The degree-one chart calculation of C.3e identifies the blow-up of the
vertex with $\operatorname{Tot}(L^{-1})$. Hence $Z$ is regular,
proper over the noetherian local ring $R$, has complement $X$, and
$p_V\circ i=1$. Its exceptional divisor has normal bundle $L^{-1}$.
The closed fibres of the $a$th infinitesimal neighbourhood of the
vertex are the split schemes
\[
 V_a=\operatorname{Spec}_V
       (\mathcal O_V\oplus L\oplus\cdots\oplus L^{a-1}).
\]
The residue field is already algebraically closed, so strict
henselization leaves the quotients by $\mathfrak m^a$ unchanged.

We apply \citet[Theorem A, equivalence (ii), equation (1.6),
author pp.~2--3]{rootKerzStrunkTamme2018} to this abstract blow-up
square. Its hypothesis is noetherianity; no resolution theorem in
positive characteristic is being assumed. The exact sequence of
pro abelian groups contains
\[
 K_0(R)\longrightarrow
 \{K_0(R/\mathfrak m_R^a)\}_a\oplus K_0(Z)
 \longrightarrow\{K_0(V_a)\}_a.
\]
All the local-ring groups in the first two positions are $\mathbb Z$,
and the first map is $n\mapsto(n,n[\mathcal O_Z])$.
Quotient by this graph. We obtain the monomorphism
$K_0(Z)\to\{K_0(V_a)\}_a$ in the abelian category of pro abelian
groups. Its kernel system consists of decreasing subgroups of the
constant group $K_0(Z)$, with inclusion transition maps. A pro-zero
system of inclusions has a zero term: for a fixed index the
eventually zero transition map is itself an inclusion. Thus some
$a_0$ satisfies
\begin{equation}\label{eq:C3ax-finite}
                   K_0(Z)\hookrightarrow K_0(V_{a_0}).
\end{equation}
This is the argument of C.3ab with the characteristic-free pro-descent
theorem in place of its characteristic-zero input.

\emph{2. Nilpotent thickenings after inverting $p$.}
For a ring $A$ of characteristic $p$ and a nilpotent ideal $I$,
\citet[Theorem A, p.~29]{rootGeisserHesselholt2011} prove that every
nonconnective relative group $K_q(A,I)$ is $p$-primary torsion
of bounded exponent. Consequently $K(A)[1/p]\to K(A/I)[1/p]$
is an equivalence. This statement globalizes to the nilpotent
immersion $V\hookrightarrow V_a$. Indeed choose a finite affine
open cover of the separated scheme $V_a$. Its intersections are
affine; on each of them the assertion just proved applies.
The Mayer--Vietoris theorem for a union of two quasi-compact open
subschemes, \citet[Theorem 8.1]{D_ThomasonTrobaugh1990}, and induction
on the finite cover give the same equivalence globally. Localization
of spectra at $p$ commutes with these finite homotopy limits.
In particular
\begin{equation}\label{eq:C3ax-nil}
                 K_0(V_a)[1/p]\xrightarrow{\sim}K_0(V)[1/p].
\end{equation}
We do not require a uniform exponent as $a$ varies, or even as
the $K$-degree varies.

Flatness of $\mathbb Z[1/p]$ preserves the injection
\eqref{eq:C3ax-finite}. Combining it with \eqref{eq:C3ax-nil}
shows that $i^*:K_0(Z)[1/p]\to K_0(V)[1/p]$ is injective.
It is surjective because $i^*p_V^*=1$. This proves
\eqref{eq:C3ax-K}.

\emph{3. Passage to rational Chow groups and ordinary families.}
The rational Chern character on a smooth quasi-projective variety
is an isomorphism from $K_0\otimes\mathbb Q$ to its full graded
rational Chow ring, naturally for pullback
\citep[Example 15.2.16(b), pp.~294--295]{frC4Fulton1998}.
The same assertion on $Z$ follows by
descent of each finite collection of vector bundles, cycles and
rational equivalences to a finite pointed-\'etale neighbourhood
of the cone vertex. These neighbourhoods have smooth
quasi-projective pullbacks to $\operatorname{Tot}(L^{-1})$;
the regular pullback and Chern-character identities there pass
to their filtered limit. The maps on Chow rings preserve codimension,
so the resulting isomorphism of direct sums from
\eqref{eq:C3ax-K} is an isomorphism in every degree. This proves
\eqref{eq:C3ax-Chow} without an integral Chern-character assertion.

Both maps in \eqref{eq:C3ax-Chow} preserve ordinary algebraic
equivalence. Pullback $p_V^*$ uses the flat smooth-limit map
of C.3e. For $i^!$, apply refined Cartier pullback to a cycle
on the actual product $Z\times C$, where $C$ is a smooth
connected projective parameter curve; its pullback is a cycle
on $V\times C$ and specialization to the two parameter points
commutes with this refined pullback. No formal parameter family
or local-Picard equivalence is used. The isomorphism therefore
also identifies $B^r(Z)_{\mathbb Q}$ with $B^r(V)_{\mathbb Q}$.

Chow localization for $V\subset Z$ has cokernel
$\mathrm{CH}^r(X)_{\mathbb Q}$. Under \eqref{eq:C3ax-Chow},
its boundary image is $-c_1(L)\mathrm{CH}^{r-1}(V)_{\mathbb Q}$,
by $i^!i_*=-c_1(L)$. This gives \eqref{eq:C3ax-open}.
Localization is also right-exact after quotienting by ordinary
algebraic equivalence: close any actual family on $X\times C$
in $Z\times C$; the difference between its fibre closures and
the prescribed fibres is supported on $V$. Equivalently one
uses the ordinary-family localization already proved in C.3e.
We obtain \eqref{eq:C3ax-B}.

Finally $B^d(V)\simeq\mathbb Z$ by degree. Any two points can
be joined by a connected projective curve on $V$, and the
diagonal on its normalization gives their ordinary algebraic
equivalence; smooth connectedness and projectivity give such
curve chains. Amplitude supplies an integral curve of positive
$L$-degree (take a complete intersection, or $V$ itself if $d=1$).
Thus multiplication by $c_1(L)$ has nonzero image in the
degree quotient, and its rational cokernel is zero.
\end{proof}
